While refrigeration cycles have been the subject of extensive efficiency studies in recent years, 
HTHP cycles have received comparatively little attention. Our objective is to implement new design 
ideas inspired by these studies, such as the dual evaporator or the transcritical cycle, to 
improve the efficiency of HTHP cycles. \\

This work is addressed to the scientific community, but also to industry, 
since it aims to provide insights and solutions that can be directly applied in real-world 
applications. In addition, a prototype will be developed to validate the proposed concepts. \\

Our work is also multidisciplinary, as it includes thermodynamics, heat transfer, fluid mechanics,
 control techniques, and machine design. It is a complete project, from the numerical design to 
 the building and testing of a prototype. \\

The research question that will be addressed in this master's thesis is the following:

%\begin{center}
%\textit{How do the integration of a dual evaporator and the operation in a transcritical 
%cycle influence the performance of a high-temperature heat pump under real-life 
%application conditions?}
%\end{center}

%\begin{center}
%\textit{How can a modular high temperature heat pump be designed, sized, 
%controlled, and prototyped to operate in both subcritical and transcritical modes, 
%with optional use of recuperators and dual evaporators?}
%\end{center}

\begin{center}
    \textit{How can a modular high-temperature heat pump be designed, 
    sized, controlled, and prototyped to operate in both subcritical and transcritical modes, 
    with optional use of recuperators and the flexibility to employ one or two evaporators?}
\end{center}